% Direct quadratic coefficient tests; preserves the earlier 114-operation milestone.
\section{Testing quadratic residuals with one mask}\label{sec:quadratic-masks}

The sum of squares used to obtain a quartic polynomial is not needed for
coefficient testing. We can test the original quadratic residuals at
separate digit positions, using the already decoded variable mask as the
test mask. Extra digit coordinates make this possible without extra
unknowns in the final polynomial system.

\subsection{A fixed layout for all quadratic systems}

Start from the quadratic-system representation in 58 existential
variables underlying the universal pair $(58,4)$ in~\cite[\S5]{J82}.
Shift its positive variables to nonnegative variables
$X_1,\ldots,X_{58}$, and put $X_0=x$. The integer quadratic residuals
have fixed coefficients for the represented r.e.\ set. Their rational
linear span has dimension at most
\[
 m=\binom{61}{2}=1830,
\]
the number of monomials of degree at most two in 59 variables.
Choose a basis from the original integer residuals and pad it with zero
polynomials to obtain $F_0,\ldots,F_{m-1}$. Simultaneous vanishing is
unchanged: every original residual is a rational linear combination of
the chosen basis, and every chosen residual is original.

Use the following fixed positions:
\begin{align*}
 v_i&=1+3(6729i+i^2)\quad(0\le i\le58),& v_*&=1180939,\\
 D&=4v_*+1,& T&=(m+1)D+2v_*,\\
 t_j&=T+jD\quad(0\le j<m),& t_*&=T+(m-1)D.
\end{align*}
Let $d_*=T-2v_*$. The inequalities
\begin{equation}\label{eq:quad-layout-bounds}
 D>2v_*,\qquad
 d_*>\max(v_*,t_*-T),\qquad
 L=5^{16}=152587890625>3t_*+1=51873937495
\end{equation}
are exact integer inequalities; the companion program checks them.
The positions $v_i$ are for the input and the 58 actual variables.
The positions $t_j$ are for 1830 additional, unconstrained digit
coordinates and will also be the equation-testing positions.

For a multiindex $I=(I_0,\ldots,I_{58})$ with $|I|\le2$, put
\[
 w(I)=\sum_{i=0}^{58}I_i v_i,
 \qquad c_I=\frac{2!}{I_0!\cdots I_{58}!(2-|I|)!}\in\{1,2\}.
\]
Write $F_j=\sum_I a_{j,I}X^I$, and define
$P_{j,I}=2a_{j,I}/c_I\in\mathbb Z$. Choose a power of two $z\ge2$
strictly larger than every $|P_{j,I}|$, and put $Z=2z$.
Define the polynomial with possibly negative coefficients
\[
 \mathcal D(B)=\sum_{j=0}^{m-1}\sum_{|I|\le2}
        P_{j,I}B^{t_j-w(I)}.
\]
Its exponents are nonnegative. Within one row, distinct multiindices
have distinct weights. Indeed $v_i\equiv1\pmod3$, so the weight
determines the total degree, which is at most two. Degree-one weights
are strictly increasing. For a degree-two weight $v_i+v_j$, subtracting
$2$ and dividing by $3$ gives $6729(i+j)+i^2+j^2$. Since
$0\le i^2+j^2\le2\cdot58^2<6729$, division by $6729$ recovers both
$i+j$ and $i^2+j^2$. These determine the unordered pair $\{i,j\}$.
The supports of different rows lie in disjoint intervals
$[t_j-2v_*,t_j]$, by \eqref{eq:quad-layout-bounds}. Thus every
coefficient of $\mathcal D$ has absolute value less than $z$, and its
support lies in $[d_*,t_*]$.

The two nonnegative digit polynomials are
\begin{align*}
 \ell_0(B)&=\sum_{i=1}^{58}B^{v_i}+\sum_{j=0}^{m-1}B^{t_j},\\
 e_0(B)&=z\sum_{h=0}^{L-1}B^h+\mathcal D(B).
\end{align*}
Every coefficient of $e_0$ lies in $[0,Z)$, and its degree is $L-1$.
The mask $\ell_0$ has coefficients zero or one and degree $t_*<L$.
The admissible index is $(Z,V,H)$, where
\begin{equation}\label{eq:quad-index}
 V=e_0(Z)+\ell_0(Z)Z^L,\qquad
 H>\max\{2Z^{2L+1},\ Z(1890)^2,\ 3L,\ 16\}
\end{equation}
and $H$ is a power of two. There are $58+1830=1888$ non-input digit
coordinates; their number explains the factor $1890=1888+2$.
They are encoded in one integer $g$ and are not new unknowns of the
final polynomial system. All these choices are effective index
construction. The certificate parameters remain $x,Z,V,H$.

\subsection{The changed equations and the bounds before decoding}

Retain the signed Pell parameter and the linear interval equations
proved above. Put $B=Hb^2$ and $\mathcal C=1+xB+g$, with both
abbreviations computed by charged instructions. Replace the bound and
the third testing block by
\begin{align*}
 e+\ell\mathcal C^2+\alpha&=q,\\
 S_3&=(2e-Z\lambda)\mathcal C^2+B\lambda(1+q^2),\\
 T_3&=(B-2)\ell.
\end{align*}
The other coding equations remain
\[
 b=x+\beta,\quad \lambda+q^2=1+\lambda B,\quad
 \theta+Z=B,\quad e+\ell q=V+t\theta,\quad n=q^8.
\]
The first two pairs of blocks are
\[
 S_1=g,\quad T_1=q-1-(b-1)\ell,\qquad
 S_2=e+\ell q,\quad T_2=\theta\lambda,
\]
with widths $N_1=q$, $N_2=q^2$, $N_3=q^5$. Accordingly,
\begin{align*}
 S&=g+q(e+\ell q+q^2S_3),\\
 T+1&=q-(b-1)\ell+\theta\lambda q+(B-2)\ell q^3,\\
 r&=S(n^2-n)+(T+1)(n^2-1).
\end{align*}

These blocks have the required signs and widths before either
$b$ being a power of two or $q=B^L$ is known. Indeed, positivity gives
$e,\ell,\mathcal C^2<q$, while $\mathcal C>B>b\ge2$ gives $q>B$.
The geometric equation yields $\lambda>q$, so
$0<z\lambda-e<z\lambda$. Moreover
\[
 2\mathcal C^2(z\lambda-e)<2z\lambda q<B\lambda q,
 \qquad 0<S_3<B\lambda(1+q^2)<2q^4<q^5.
\]
Since $(b-1)\ell<\ell\mathcal C^2<q$, we have $0\le T_1<q$.
Likewise $0<T_3=(B-2)\ell<\ell\mathcal C^2<q$.
The familiar integer bounds give $0<S_2<q^2$ and
$0<T_2<(B-1)\lambda=q^2-1$, while $0<g<q$.
Consequently $r\ge n^2-1\ge n$, and admissibility gives
\[
 3<3L\le B<q\le n\le r,
 \qquad r,n\ge8,\qquad 0<b\le n.
\]
As $B$ is even already, the noncircular Pell and interval arguments of
Section~\ref{sec:further-pell} apply. They give $b$ a power of two,
$q=B^L$, and $n^2\mid\binom{2r}{r}$. Thus the three separate no-carry
conditions follow from Lemmas~2.16 and~2.11 of~\cite{J82}.

\subsection{Decoding and excluding unwanted coefficients}

Apply Lemma~2.9 of~\cite{J82} to $Y=e+\ell q$, with digit base $Z$
and $m=k=2L$. Its fixed digit polynomial is
$e_0(B)+\ell_0(B)B^L$, of degree $L+t_*<2L$, with coefficients
in $[0,Z)$. Equation~\eqref{eq:quad-index} supplies its value at $Z$;
the packed congruence, $Y<q^2=B^{2L}$, the mask $\theta\lambda$, and
$B>2Z^{2L+1}$ supply the other hypotheses. Therefore
$e+\ell q=e_0(B)+\ell_0(B)q$. Since all four components lie in
$[0,q)$, quotient-and-remainder uniqueness gives
$e=e_0(B)$ and $\ell=\ell_0(B)$.

The first no-carry condition now decodes
\[
 \mathcal C(B)=1+xB+
       \sum_{i=1}^{58}X_i B^{v_i}
       +\sum_{j=0}^{m-1}Y_j B^{t_j},
 \qquad 0\le X_i,Y_j<b.
\]
The $Y_j$ are the unconstrained extra coordinates. They do not affect
the equation tests. To see this, regard $B$ as a formal variable and
write
\[
 e_0-z\lambda
   =\mathcal D-z\sum_{h=L}^{2L-1}B^h.
\]
The tail starts above every testing position, since $t_*<L$.
The remaining factor $\mathcal D$ starts at degree $d_*$.
Any monomial of $\mathcal C^2$ involving a $Y_j$ has degree at least
$T$, so its product with $\mathcal D$ has degree at least
$T+d_*>t_*$. Thus no extra coordinate contributes to any testing
position. Also $d_*>v_*$, so the coefficients at the low positions
$v_i$ are automatically zero.

At a high position $t_j$, a contribution from row $k$ and true-variable
monomials $I,J$ requires
\[
 t_k-w(I)+w(J)=t_j.
\]
Both weights lie in $[0,2v_*]$, so the row spacing forces $k=j$.
Uniqueness of the quadratic monomial weights then forces $I=J$. Therefore
\begin{equation}\label{eq:quad-target-coefficient}
 [B^{t_j}]\,\mathcal C(B)^2(e_0(B)-z\lambda(B))
       =\sum_I c_I P_{j,I}X^I=2F_j(x,X_1,\ldots,X_{58}).
\end{equation}
This proves both kinds of separation: other rows cannot interfere,
and the extra coordinates at the row positions cannot interfere.

Every coefficient of $e_0-z\lambda$ has absolute value at most $z$.
The sum of coefficients of $\mathcal C^2$ is less than
$(1890b)^2$. Hence every coefficient of their product has absolute
value less than
\[
 z(1890b)^2<B/2.
\]
Its degree is at most $2L-1+2t_*<3L$. The offset
$(B/2)\lambda(1+q^2)$ supplies the digit $B/2$ at every position
below $4L$, so all resulting digits lie strictly between zero and $B$.
Now $S_3,T_3$ are even, and their third no-carry condition is equivalent
to that for
\[
 \frac{S_3}{2}\quad\hbox{and}\quad
 \frac{T_3}{2}=(B/2-1)\ell_0(B).
\]
The latter mask tests exactly the low positions $v_i$ and the high
positions $t_j$. At each such position it requires the digit of
$S_3/2$ to equal $B/2$. The low tests hold automatically; by
\eqref{eq:quad-target-coefficient}, the high tests are equivalent to
$F_j=0$ for every $j$. This proves soundness.

\subsection{Positive witnesses and the certificate count}

Conversely, choose a nonnegative solution of the quadratic system and
assign the extra coordinates arbitrarily in $[0,b)$, with one of them
equal to $1$. Choose the power of two $b$ larger than $x$ and all these
coordinates. The extra value $1$ guarantees $g>0$ even if every true
existential coordinate is zero. Take the intended $B,q,e,\ell,g$.
Their bounds are
\[
 e<ZB^{L-1}<q/2,\qquad
 \ell<2B^{t_*},\qquad
 \mathcal C<2bB^{t_*}.
\]
Using $H>16$ and $L>3t_*+1$, we obtain
\[
 \ell\mathcal C^2<8b^2B^{3t_*}
      <\tfrac12B^{3t_*+1}\le q/2.
\]
Thus $\alpha=q-e-\ell\mathcal C^2$ is positive, as is $\beta=b-x$.
The packed quotient $t$ is a positive integer because its fixed
nonnegative digit polynomial has positive degree and $B>Z$.
All decoded no-carry conditions hold by the preceding coefficient
calculation. The remaining positive witnesses follow from the Pell
necessity constructions and the positive interval proof. This proves
necessity with all domains preserved.

\begin{theorem}\label{thm:best113}
For each r.e.\ set, an effective admissible index $(Z,V,H)$ as above
gives a system in 33 positive unknowns and 21 equations such that, for
positive input $x$, membership is equivalent to a valid certificate with
\textbf{113 arithmetic instructions}: 62 multiplications and 51 additions.
Equality tests and supplied parameters are free. Generating the fixed
numerals $2,4,5,L$ from $1$ adds seven instructions, giving 120 under
that convention.
\end{theorem}
\begin{proof}
Relative to the 114-instruction certificate, $B=Hb^2$ saves the
multiplication constructing $b^4$, and $\mathcal C^2$ saves the one
constructing $\mathcal C^4$. Computing the changed third mask in the
packing expression costs one extra multiplication; the other counts
are unchanged. The $T+1$ computation may be organized as
\[
 T+1=q(\lambda+q^2-Z\lambda)
      +(B-2)\bigl((\ell q)q^2\bigr)-(b-1)\ell,
\]
using the geometric and $\theta$ equations and the already computed
product $\ell q$. The exact schedule and
its triangular residual corrections are verified by
\file{round6\_1980\_quadratic\_certificate.py}.
The source system has the same 33 positive unknowns and 21 equations;
the additional digit coordinates are encoded inside $g$.
The equivalence is proved above. Starting from $1$, three additions
produce $2,4,5$ and four successive squarings of $5$ produce $L=5^{16}$,
giving the seven-instruction numeral construction.
\end{proof}
